\section{An exact finite indicator-growth bridge and its limits}
\label{sec:p9-indicator-bridge}

This section provides a deliberately restricted population model in which
function-preserving growth and the deterministic partition hierarchy problem
coincide. It is a specialization of proper-loss regret and the Bregman
centroid identity, not a general neural-training result or a priority claim.

Fix an integer $N\ge2$, masses $p_i>0$ with $\sum_{i=1}^N p_i=1$, and
posterior values $x_i\in[0,1]$. The admissible inputs are the standard basis
vectors $e_1,\ldots,e_N\in\mathbb R^N$, and the population law is
\[
 \Pr(X=e_i)=p_i,\qquad
 \Pr(Y=1\mid X=e_i)=x_i.
\]
Let $\varphi\in C^1([0,1])$ be strictly convex. For definiteness use the
proper loss
\[
 \ell_\varphi(q,y)=-\varphi(q)-\varphi'(q)(y-q),
 \quad q\in[0,1],\quad y\in\{0,1\}.
\]
Adding a constant depending only on $y$ changes none of the excess risks.
Define $B_\varphi(u,v)=\varphi(u)-\varphi(v)-\varphi'(v)(u-v)$.
The conditional regret of report $q$ at posterior $x$ is exactly
$B_\varphi(x,q)$, so strict convexity makes the loss strictly proper.

For a partition $P$ of $\{1,\ldots,N\}$, let $w_C\in\{0,1\}^N$ be
the indicator of $C\in P$, and use the restricted ReLU network
\[
 f_{P,q}(z)=\sum_{C\in P}q_C\operatorname{ReLU}(w_C^\top z),
 \qquad q_C\in[0,1].
\]
On every admissible input $e_i$, this equals the report $q_C$ of its
unique cell. Capacity counts the active leaf units $|P|$; there are no
additional depth, split-search, sample, parameter-storage, or training-cost
constraints. Hidden weights are constrained to these disjoint indicator
vectors. Their ambient dimension $N$ is not charged by the leaf count.
For $C\in P$, put $p_C=\sum_{i\in C}p_i$ and
$m_C=p_C^{-1}\sum_{i\in C}p_ix_i$, and define
\[
 D_\varphi(P)=\sum_{C\in P}\sum_{i\in C}
 p_i\bigl(\varphi(x_i)-\varphi(m_C)\bigr),\qquad
 \operatorname{OPT}_k=\min_{|P|\le k}D_\varphi(P).
\]
Write $R(P,q)=\mathbb E\ell_\varphi(f_{P,q}(X),Y)$ and
$R^*=\mathbb E\ell_\varphi(x_X,Y)$, with $x_{e_i}=x_i$.

\begin{theorem}[Exact oracle equivalence for indicator splitting]
\label{thm:p9-bridge}
For every finite population just specified, the following statements hold.
For every partition $P$ and every report vector $q$,
\begin{equation}
 R(P,q)-R^*=D_\varphi(P)+
       \sum_{C\in P}p_CB_\varphi(m_C,q_C).
 \label{eq:p9-risk-decomposition}
\end{equation}
Thus the unique population-optimal reports are $q_C=m_C$, and the resulting
excess risk is $D_\varphi(P)$.

Splitting a cell $C=C_0\sqcup C_1$ into two nonempty children and copying
its report into both children preserves the current predictor on every
admissible input. If arbitrary such subset splits are permitted, and exact
population report minimization follows each birth, the smallest simultaneous
risk factor of a one-leaf-to-$N$-leaf growth path is exactly
\begin{equation}
 \rho_{\rm ind}
 =\min_{\substack{H_{k+1}\text{ refines }H_k\\|H_k|\le k}}
 \inf\{c\ge0:D_\varphi(H_k)\le c\operatorname{OPT}_k
                       \text{ for every }1\le k\le N\}.
 \label{eq:p9-price-equivalence}
\end{equation}
Here the growth-path factor uses $R_k-R^*$ in place of $D_\varphi(H_k)$.
The cost-inequality convention includes capacities with zero flat optimum.
The minimum is over complete schedules selected with knowledge of the
population; it does not assert a greedy or sample-based selection rule.
\end{theorem}
\begin{proof}
For each cell, expand and add the Bregman divergences:
\[
 \sum_{i\in C}p_iB_\varphi(x_i,q_C)
 =\sum_{i\in C}p_i\bigl(\varphi(x_i)-\varphi(m_C)\bigr)
  +p_CB_\varphi(m_C,q_C).
\]
The linear term involving $x_i-m_C$ cancels. Summing proves
\eqref{eq:p9-risk-decomposition}; strict convexity proves uniqueness.
The same identity holds conditional on any realized learned partition and
reports, provided evaluation uses a fresh independent draw from the fixed
population law above.

At a split, for all $i$ we have
$\mathbf1_C(i)=\mathbf1_{C_0}(i)+\mathbf1_{C_1}(i)$. Equal copied
reports therefore preserve $f(e_i)$, and the old restricted function class
is contained in the new one. In this implementation the same equality
also holds on the nonnegative orthant; no claim of preservation on all of
$\mathbb R^N$ is needed. Exact population refitting then attains
$D_\varphi(P')$ for the refined partition $P'$.

Every binary split path is an admissible hierarchy, giving one inequality
in \eqref{eq:p9-price-equivalence}. For the converse, take any admissible
hierarchy $(H_k)$, which may have fewer than $k$ cells. Let $Q_N$ be the
singleton partition. Induct backwards. If $Q_{k+1}$ has $k+1$ cells and
refines $H_{k+1}$, it also refines $H_k$. Since $|H_k|\le k$, two of its
cells lie inside the same $H_k$ cell. Merge these two to form $Q_k$.
Then $Q_k$ has exactly $k$ cells and refines $H_k$. Refinement and
convexity give $D_\varphi(Q_k)\le D_\varphi(H_k)$. Reversing this
construction yields an admissible binary split path with no larger cost
at any capacity. Oracle refitting realizes all its costs. This proves the
reverse inequality, including zero optima.
\end{proof}

If the reports before and after a split are population optimal, its gain is
\[
 D_\varphi(P)-D_\varphi(P')
 =\sum_{j=0}^1p_{C_j}B_\varphi(m_{C_j},m_C)\ge0.
\]
The gain is zero when the two child means coincide. A realized empirical
refit instead retains the explicit nonnegative report-error term in
\eqref{eq:p9-risk-decomposition}. The theorem is an equivalence of
population oracle problems, not a convergence or generalization theorem.

\begin{proposition}[Neuron count and birth alone do not transfer the result]
\label{prop:p9-nontransfer}
On the same finite input domain and population law, removing the indicator
constraint on hidden weights allows the Bayes predictor to be represented
by a single ReLU unit. Also, even in the restricted indicator model,
function-preserving birth alone need not achieve any finite factor relative
to the flat optimum at the new capacity.
\end{proposition}
\begin{proof}
Set $w=(x_1,\ldots,x_N)$ and $f(z)=\operatorname{ReLU}(w^\top z)$.
Then $f(e_i)=x_i$ because $x_i\ge0$, so its population excess risk is zero.
This uses one hidden unit with output coefficient one. Therefore replacing
leaf count by the width of this unrestricted ReLU class changes the flat
oracle problem completely. In particular, high indicator-hierarchy prices
cannot imply a width lower bound for this larger network class.

For the second claim take $N=2$, equal masses, $x_1=1/4$, $x_2=3/4$, and
squared loss $(y-q)^2$. The optimal one-leaf report is $1/2$ and its excess
risk is $1/16$. After a singleton split that copies report $1/2$, the
function and excess risk are unchanged, whereas $\operatorname{OPT}_2=0$.
Thus this birth, without a refit guarantee, violates every finite
multiplicative cost inequality at capacity two.
\end{proof}

\paragraph{Absolute normalization.}
For every $\lambda>0$, replacing the loss and generator by
$\lambda\ell$ and $\lambda\varphi$ scales every excess risk and every
$\operatorname{OPT}_k$ by $\lambda$, preserving all hierarchy prices.
If a loss already lies in $[0,1/2]$, every $0<\lambda\le1$ retains that
same bounded-loss requirement. Given any finite instance and any
$\epsilon>0$, taking $\lambda$ sufficiently small makes all its
population excess risks below $\epsilon$ without changing its price.
Consequently, a divergent normalized relative price alone cannot imply
a nonvanishing absolute-risk separation under this upper-bound
normalization. An exact fixed-scale normalization would need its own
analysis.

A stochastic degradation chain is a different representation class: the
deterministic indicator theorem does not identify its price with that of
physical randomized neuron growth. The theorem requires the population
posterior oracle and exact refits; it supplies no empirical memory,
split-search, sample, or optimizer guarantee.
